% ============================================================
%  LEMBAR SOAL  --  MATEMATIKA  --  SSU-ITB 2026  |  Paket 5
%  Sumber: SM-SSU ITB 2025 (50 soal)
% ============================================================

\documentclass[11pt]{article}

\usepackage[utf8]{inputenc}
\usepackage[T1]{fontenc}
\usepackage[a4paper,margin=2.5cm,top=2.8cm]{geometry}
\usepackage{amsmath,amssymb}
\usepackage{graphicx}
\usepackage{enumitem}
\usepackage{multicol}
\usepackage{xcolor}
\usepackage[most]{tcolorbox}
\usepackage{fancyhdr}
\usepackage{booktabs}
\usepackage{icomma}
\usepackage{array}

\definecolor{mapelcolor}{RGB}{20,110,60}
\definecolor{mapelbg}{RGB}{228,246,234}

\pagestyle{fancy}\fancyhf{}
\lhead{\small SSU-ITB 2026}
\rhead{\small Paket 5}
\cfoot{\small\thepage}
\renewcommand{\headrulewidth}{0.4pt}

\newtcolorbox{soal}[1]{
  breakable,
  colback=white, colframe=black!70,
  coltitle=white, colbacktitle=black,
  fonttitle=\bfseries\sffamily,
  title=Nomor #1,
  boxrule=0.8pt, arc=3pt,
  left=3mm, right=3mm, top=2mm, bottom=2mm}

\newenvironment{pil}{%
  \begin{enumerate}[label=(\alph*),leftmargin=*,itemsep=1pt,topsep=3pt]}%
  {\end{enumerate}}
\newenvironment{pilB}{%
  \par\smallskip\begin{multicols}{2}
  \begin{enumerate}[label=(\alph*),leftmargin=*,itemsep=1pt,topsep=3pt]}%
  {\end{enumerate}\end{multicols}}

\newcommand{\gambar}[2][0.45\textwidth]{%
  \begin{center}\includegraphics[width=#1]{#2}\end{center}}

\linespread{1.05}

% ===================================================================
\begin{document}
\thispagestyle{empty}

\begin{center}
  \fbox{\parbox{0.6\textwidth}{\centering\bfseries\large LEMBAR SOAL}}
\end{center}
\vspace{0.5cm}
\begin{center}{\LARGE\bfseries MATEMATIKA}\end{center}
\vspace{0.5cm}

\begin{center}
  \begin{tabular}{@{\hspace{2cm}}ll@{\hspace{0.4cm}:@{\hspace{0.4cm}}}l}
    Jenjang      & & SMA / Sederajat               \\[0.3em]
    Hari/Tanggal & & \underline{\hspace{5cm}}      \\[0.3em]
    Waktu        & & \underline{\hspace{5cm}}      \\[0.3em]
    Paket Soal   & & 5                             \\
  \end{tabular}
\end{center}

\vspace{0.6cm}\noindent\rule{\linewidth}{0.6pt}\vspace{0.3cm}

\noindent\textbf{\underline{PETUNJUK UMUM}}
\begin{enumerate}[leftmargin=*,itemsep=4pt,topsep=4pt]
  \item Anda \textbf{tidak} diperbolehkan menggunakan sumber-sumber eksternal,
        termasuk internet, \textit{large language model} seperti ChatGPT, Claude,
        LLaMA, Gemini, Meta AI, dan sebagainya. Anda sangat dianjurkan untuk
        mencoba mengerjakan sendiri terlebih dahulu karena evaluasi ini dibuat
        untuk \textbf{mengukur} tingkat kepahaman; nilai $<70$ mengindikasikan
        \textbf{tidak lulus}, dan jika sudah melewati \textit{deadline} maka
        dianggap tidak mengerjakan yang berarti \textbf{tidak lulus}.

  \item Terdapat 2 tahap pengerjaan, yaitu tahap \textbf{Try Out} dan tahap
        \textbf{Review}. Try Out bersifat \textit{open book}; Anda diperbolehkan
        membuka buku apapun selama bersifat \textit{offline} dengan rentang waktu
        2 jam pada tanggal \underline{\hspace{3cm}}. Setelah itu memasuki
        tahap Review dengan \textit{schedule} rentang tanggal
        \underline{\hspace{3cm}}--\underline{\hspace{3cm}}, di mana Anda
        diharapkan mengerjakan ulang dengan disertakan penjelasan/pembelaan
        mengapa revisi jawaban adalah pilihan tersebut, dan tetap mencantumkan
        penjelasan pada soal yang walaupun pada penilaian sudah benar.

  \item Anda dianjurkan untuk mencantumkan caranya walaupun tidak termasuk
        penilaian, mengingat sifatnya adalah sebagai bahan pembelajaran.
        Mohon bertanggung jawab atas pekerjaan Anda sendiri.

  \item Bacalah setiap soal dengan teliti. Terdapat tiga jenis soal:
        \begin{enumerate}[label=(\arabic*),leftmargin=2em,topsep=2pt,itemsep=2pt]
          \item \textbf{Pilihan Ganda Biasa} --- 5 opsi (A--E), pilih
                1 jawaban paling tepat;
          \item \textbf{Pilihan Ganda Majemuk Kompleks} --- tabel
                \textbf{Benar}/\textbf{Salah}, tentukan kebenaran tiap
                pernyataan secara mandiri; dan
          \item \textbf{Isian Singkat} --- tulis nilai/ekspresi pada
                kolom yang tersedia.
        \end{enumerate}

  \item Soal berjumlah \textbf{50 butir} soal dengan bobot asli per soal adalah
        \textbf{2} untuk mendapatkan nilai sempurna sebesar \textbf{100}.

  \item Silakan bertanya apabila terdapat soal yang ambigu, rancu, atau kurang
        spesifik selama itu tidak menjurus kepada jawaban soal tersebut.

  \item Isilah detail informasi berikut.\\[0.3em]
        \textbf{Nama}\hspace{0.45cm}:\enspace\underline{\hspace{8cm}}
\end{enumerate}

\vspace{0.6cm}
\noindent\textit{Dengan menandatangani kolom di bawah ini, saya menyatakan telah membaca,
memahami, dan menyetujui seluruh peraturan yang tercantum di atas.}

\vspace{0.6cm}
\noindent\fbox{\parbox{5cm}{\vspace{2.2cm}\centering\small Tanda Tangan Peserta\vspace{0.3cm}}}

\clearpage

% ===================================================================
%  SOAL 1--10
% ===================================================================

\begin{soal}{1 \normalfont\small[PG Biasa]}
Diketahui $a$ adalah bilangan real. Manakah pernyataan yang \textbf{selalu} benar?
\begin{pil}
  \item Jika $a\geq0$: $\dfrac{a}{a+1}>\dfrac{a+1}{a+2}$.
  \item Jika $a<0$ dan $a\neq-1$: $\dfrac{1}{a+1}<\dfrac{1}{a}$.
  \item Jika $a<0$: $2+a<1-a$.
  \item Jika $a>0$: $a>1-a$.
  \item Jika $a>0$: $-\dfrac{1}{a}<-\dfrac{1}{a+1}$.
\end{pil}
\end{soal}

\begin{soal}{2 \normalfont\small[PG Majemuk Kompleks]}
Diberikan $f(x)=\dfrac{1}{\sqrt{2x^2-5x-3}}$.
Tentukan kebenaran setiap pernyataan berikut!

\par\smallskip
{\renewcommand{\arraystretch}{1.4}%
\begin{tabular}{@{}p{0.76\linewidth}cc@{}}
\hline\textbf{Pernyataan}&\textbf{Benar}&\textbf{Salah}\\\hline
(1)~Syarat terdefinisi: $2x^2-5x-3>0$.&$\bigcirc$&$\bigcirc$\\
(2)~Faktorkan: $2x^2-5x-3=(2x+1)(x-3)>0$, terpenuhi untuk $x<-\tfrac12$ atau $x>3$.&$\bigcirc$&$\bigcirc$\\
(3)~$x=4$ memenuhi syarat terdefinisi.&$\bigcirc$&$\bigcirc$\\
(4)~$x=5$ memenuhi syarat terdefinisi.&$\bigcirc$&$\bigcirc$\\\hline
\end{tabular}}
\end{soal}

\begin{soal}{3 \normalfont\small[Isian Singkat]}
Di suatu pusat perbelanjaan, peluang seseorang makan, membeli buku, dan
membeli mainan masing-masing $\frac{2}{5}$, $\frac{3}{7}$, $\frac{7}{8}$
(saling bebas). Peluang membeli buku atau mainan tetapi tidak makan
adalah $\dfrac{P}{840}$. Nilai $P=\ldots$

\vspace{0.4em}\noindent Jawaban:\enspace\underline{\hspace{4cm}}
\end{soal}

\begin{soal}{4 \normalfont\small[Isian Singkat]}
Titik-titik potong kurva $y=(x+3)^2$ dengan garis $y=px+9$ adalah
$(-3,0)$ dan $(r,s)$. Nilai $p+r+s=\ldots$

\vspace{0.4em}\noindent Jawaban:\enspace\underline{\hspace{4cm}}
\end{soal}

\begin{soal}{5 \normalfont\small[PG Biasa]}
\gambar[0.24\textwidth]{images/5/5-2.png}
Kerucut tegak memiliki tinggi $\sqrt{77}$ cm dan diameter 4 cm.
Panjang garis pelukis $s$ adalah \ldots\ cm.
\begin{pilB}
  \item $9$ \item $11$ \item $13$ \item $8$ \item $10$
\end{pilB}
\end{soal}

\begin{soal}{6 \normalfont\small[PG Biasa]}
Jika $y=\dfrac{3x+7}{x+9}$, maka $x=\ldots$
\begin{pil}
  \item $\dfrac{7-9y}{y-3}$ \item $\dfrac{7+9y}{y+3}$
  \item $\dfrac{y+3}{9+7y}$ \item $\dfrac{y+3}{9-7y}$
  \item $\dfrac{9-7y}{y-3}$
\end{pil}
\end{soal}

\begin{soal}{7 \normalfont\small[PG Biasa]}
Sistem persamaan yang mempunyai solusi $(x,y)=(2,3)$ adalah \ldots
\begin{pil}
  \item $2x+y=7$;\enspace$x-2y=-4$
  \item $2x+y=6$;\enspace$x-2y=-4$
  \item $2x+y=-4$;\enspace$x-2y=7$
  \item $2x+y=6$;\enspace$x-2y=-3$
  \item $2x+y=7$;\enspace$x-2y=-3$
\end{pil}
\end{soal}

\begin{soal}{8 \normalfont\small[Isian Singkat]}
Misalkan $f(x)=a^{x+b}$ untuk konstanta $a$ dan $b$. Jika $f(1)=3$ dan
$f(3)=18$, nilai $a^4=\ldots$

\vspace{0.4em}\noindent Jawaban:\enspace\underline{\hspace{4cm}}
\end{soal}

\begin{soal}{9 \normalfont\small[Isian Singkat]}
Garis $l$ bergradien $\dfrac{1}{5}$ dan melalui $P(7,3)$ dan $Q(k,5)$.
Nilai $k=\ldots$

\vspace{0.4em}\noindent Jawaban:\enspace\underline{\hspace{4cm}}
\end{soal}

\begin{soal}{10 \normalfont\small[PG Biasa]}
\gambar[0.4\textwidth]{images/5/5-4.png}
$AB=13$, $CD=5$, $EF=10$; ketiganya sejajar. Nilai $AE:ED=\ldots$
\begin{pilB}
  \item $3:5$ \item $13:10$ \item $8:5$ \item $13:5$ \item $10:5$
\end{pilB}
\end{soal}

% ===================================================================
%  SOAL 11--20
% ===================================================================

\begin{soal}{11 \normalfont\small[PG Biasa]}
\gambar[0.5\textwidth]{images/5/5-5.png}
Luas satu persegi kecil adalah 5. Total luas daerah yang diarsir adalah \ldots
\begin{pil}
  \item $200+\dfrac{125}{4}\pi$
  \item $20+\dfrac{15}{2}\pi$
  \item $200+\dfrac{75}{2}\pi$
  \item $40+\dfrac{25}{4}\pi$
  \item $40+\dfrac{15}{2}\pi$
\end{pil}
\end{soal}

\begin{soal}{12 \normalfont\small[PG Majemuk Kompleks]}
Diketahui $\displaystyle\int x^{2/6}\,dx=ax^b+C$ dan $a-b=\dfrac{p}{48}$.
Tentukan kebenaran setiap pernyataan berikut!

\par\smallskip
{\renewcommand{\arraystretch}{1.4}%
\begin{tabular}{@{}p{0.76\linewidth}cc@{}}
\hline\textbf{Pernyataan}&\textbf{Benar}&\textbf{Salah}\\\hline
(1)~$x^{2/6}=x^{1/3}$.&$\bigcirc$&$\bigcirc$\\
(2)~$\displaystyle\int x^{1/3}\,dx=\dfrac{3}{4}x^{4/3}+C$.&$\bigcirc$&$\bigcirc$\\
(3)~Dari $ax^b+C$: $a=\dfrac{4}{3}$ dan $b=\dfrac{3}{4}$.&$\bigcirc$&$\bigcirc$\\
(4)~Nilai $p=-28$.&$\bigcirc$&$\bigcirc$\\\hline
\end{tabular}}
\end{soal}

\begin{soal}{13 \normalfont\small[Isian Singkat]}
Data berisi 20 bilangan dengan rata-rata 42. Jika $n$ bilangan dengan
rata-rata 35 ditambahkan, rata-rata data baru menjadi 40. Nilai $n=\ldots$

\vspace{0.4em}\noindent Jawaban:\enspace\underline{\hspace{4cm}}
\end{soal}

\begin{soal}{14 \normalfont\small[PG Biasa]}
\gambar[0.4\textwidth]{images/5/5-6.png}
Diberikan grafik $y=f(x)$ seperti gambar. Jika $g(x)=f(x-2)+12$,
maka $g(4)=\ldots$
\begin{pilB}
  \item $15$ \item $12$ \item $16$ \item $13$ \item $14$
\end{pilB}
\end{soal}

\begin{soal}{15 \normalfont\small[Isian Singkat]}
Jumlah 50 suku pertama barisan $a_n=4n+7$, $n=1,2,3,\ldots$ adalah \ldots

\vspace{0.4em}\noindent Jawaban:\enspace\underline{\hspace{4cm}}
\end{soal}

\begin{soal}{16 \normalfont\small[PG Biasa]}
Jika $f(x)=\dfrac{8}{x}$, maka $\dfrac{f(7+h)-f(7)}{h}=\ldots$
\begin{pilB}
  \item $\dfrac{8}{7h-49}$ \item $\dfrac{8}{7h+49}$
  \item $-\dfrac{8}{7h+49}$ \item $\dfrac{8}{-7h+49}$
  \item $\dfrac{8}{7h+7}$
\end{pilB}
\end{soal}

\begin{soal}{17 \normalfont\small[PG Majemuk Kompleks]}
$\dfrac{1}{x+1}+\dfrac{8}{x^2+3}=\dfrac{x^2+Fx+G}{x^3+Cx^2+Dx+E}$.
Tentukan kebenaran setiap pernyataan berikut!

\par\smallskip
{\renewcommand{\arraystretch}{1.4}%
\begin{tabular}{@{}p{0.76\linewidth}cc@{}}
\hline\textbf{Pernyataan}&\textbf{Benar}&\textbf{Salah}\\\hline
(1)~Penyebut persekutuan $(x+1)(x^2+3)=x^3+x^2+4x+3$.&$\bigcirc$&$\bigcirc$\\
(2)~Pembilang gabungan: $x^2+3+8(x+1)=x^2+8x+11$.&$\bigcirc$&$\bigcirc$\\
(3)~$F=8$, $G=11$, $C=1$, $D=3$, $E=3$.&$\bigcirc$&$\bigcirc$\\
(4)~$C+D+E+F+G=26$.&$\bigcirc$&$\bigcirc$\\\hline
\end{tabular}}
\end{soal}

\begin{soal}{18 \normalfont\small[PG Biasa]}
Jika $(x+y)^2=50$ dan $xy=6$, maka $x^2+y^2=\ldots$
\begin{pilB}
  \item $62$ \item $38$ \item $50$ \item $44$ \item $56$
\end{pilB}
\end{soal}

\begin{soal}{19 \normalfont\small[PG Biasa]}
Himpunan penyelesaian $\dfrac{5}{2}\leq\dfrac{6-x}{2}<6$ adalah \ldots
\begin{pil}
  \item $\{x\mid-1<x\leq6\}$
  \item $\{x\mid-6<x\leq1\}$
  \item $\{x\mid-1\leq x<6\}$
  \item $\{x\mid-6\leq x<1\}$
  \item $\{x\mid1\leq x<6\}$
\end{pil}
\end{soal}

\begin{soal}{20 \normalfont\small[PG Biasa]}
\gambar[0.4\textwidth]{images/5/5-7.png}
Jajargenjang $ABCD$, keliling 34, luas 33, $BC=6$.
Jarak titik $D$ ke garis $AB$ adalah \ldots
\begin{pilB}
  \item $11$ \item $12$ \item $6$ \item $9$ \item $3$
\end{pilB}
\end{soal}

% ===================================================================
%  SOAL 21--30
% ===================================================================

\begin{soal}{21 \normalfont\small[PG Biasa]}
Persamaan garis singgung $y=5(x-1)^8$ pada $x=2$ adalah \ldots
\begin{pil}
  \item $y+80=40x$
  \item $y+80=5x$
  \item Tidak ada yang benar
  \item $y+75=5x$
  \item $y+75=40x$
\end{pil}
\end{soal}

\begin{soal}{22 \normalfont\small[Isian Singkat]}
Misalkan $f(x)=2x\cdot g(x)$. Jika $g(2)=2$ dan $g'(2)=6$, maka $f'(2)=\ldots$

\vspace{0.4em}\noindent Jawaban:\enspace\underline{\hspace{4cm}}
\end{soal}

\begin{soal}{23 \normalfont\small[PG Biasa]}
Jika $z:y=2:3$ dan $z^2=\dfrac{1}{6}$, maka $y^2=\ldots$
\begin{pilB}
  \item $\dfrac{3}{12}$ \item $\dfrac{2}{18}$ \item $\dfrac{1}{6}$
  \item $\dfrac{4}{54}$ \item $\dfrac{9}{24}$
\end{pilB}
\end{soal}

\begin{soal}{24 \normalfont\small[PG Majemuk Kompleks]}
Tentukan kebenaran setiap pernyataan tentang garis $x+3y=12$!

\par\smallskip
{\renewcommand{\arraystretch}{1.4}%
\begin{tabular}{@{}p{0.76\linewidth}cc@{}}
\hline\textbf{Pernyataan}&\textbf{Benar}&\textbf{Salah}\\\hline
(1)~$(6,1)$ tidak dilalui garis $x+3y=12$ karena $6+3(1)=9\neq12$.&$\bigcirc$&$\bigcirc$\\
(2)~$(12,0)$ tidak dilalui garis $x+3y=12$.&$\bigcirc$&$\bigcirc$\\
(3)~$(-3,5)$ dilalui garis karena $-3+3(5)=12$.&$\bigcirc$&$\bigcirc$\\
(4)~$(3,3)$ dilalui garis karena $3+3(3)=12$.&$\bigcirc$&$\bigcirc$\\\hline
\end{tabular}}
\end{soal}

\begin{soal}{25 \normalfont\small[Isian Singkat]}
\gambar[0.5\textwidth]{images/5/5-9.png}
$p(x)$ bernilai positif; grafik $y=p'(x)$ seperti gambar. Interval terbesar
agar $f(x)=\sqrt{p(x)+3x+1}$ turun adalah $(k,l)$. Nilai $l-k=\ldots$

\vspace{0.4em}\noindent Jawaban:\enspace\underline{\hspace{4cm}}
\end{soal}

\begin{soal}{26 \normalfont\small[Isian Singkat]}
Jika $a>0$ dan $\displaystyle\lim_{x\to\pi/a}\dfrac{1-\cos(ax-\pi)}{\left(x-\frac{\pi}{a}\right)^2}=8$,
maka $a^2+a=\ldots$

\vspace{0.4em}\noindent Jawaban:\enspace\underline{\hspace{4cm}}
\end{soal}

\begin{soal}{27 \normalfont\small[PG Biasa]}
Diketahui $f(x)=\log_2 x$ untuk $x>0$. Nilai $2f(4x)-f\!\left(\dfrac{2}{x^2}\right)=\ldots$
\begin{pilB}
  \item $f(8)+f(x^4)$
  \item $f(8)-f(x^4)$
  \item $2f(2)+f(x^4)$
  \item $f(4)\times f(x^2)$
  \item $f(2)+f(x^3)$
\end{pilB}
\end{soal}

\begin{soal}{28 \normalfont\small[PG Biasa]}
\gambar[0.3\textwidth]{images/5/5-10.png}
Segitiga $ABC$ siku-siku di $B$, $AB=BC$, $AC=8$. Luas segitiga $ABC$ adalah \ldots
\begin{pilB}
  \item $18$ \item $22$ \item $16$ \item $20$ \item $32$
\end{pilB}
\end{soal}

\begin{soal}{29 \normalfont\small[PG Biasa]}
$\left|\dfrac{x}{3}-9\right|\leq1$ menghasilkan $A\leq x\leq B$. Nilai $A+B=\ldots$
\begin{pilB}
  \item $54$ \item $30$ \item $-18$ \item $0$ \item $78$
\end{pilB}
\end{soal}

\begin{soal}{30 \normalfont\small[Isian Singkat]}
Jumlah semua nilai $n$ yang memenuhi $\log_9(89-n^2)=\log_3(n+3)$ adalah \ldots

\vspace{0.4em}\noindent Jawaban:\enspace\underline{\hspace{4cm}}
\end{soal}

% ===================================================================
%  SOAL 31--40
% ===================================================================

\begin{soal}{31 \normalfont\small[PG Majemuk Kompleks]}
Kurva $y=2x^3-3x^2+11$; tentukan garis singgung di $P(-1,q)$.
Tentukan kebenaran setiap pernyataan berikut!

\par\smallskip
{\renewcommand{\arraystretch}{1.4}%
\begin{tabular}{@{}p{0.76\linewidth}cc@{}}
\hline\textbf{Pernyataan}&\textbf{Benar}&\textbf{Salah}\\\hline
(1)~$q = f(-1) = 2(-1)^3-3(-1)^2+11 = 6$.&$\bigcirc$&$\bigcirc$\\
(2)~Gradien garis singgung $m = y'(-1) = 6(1)+6 = 12$.&$\bigcirc$&$\bigcirc$\\
(3)~Persamaan garis singgung di $(-1,6)$ adalah $y=12x+12$.&$\bigcirc$&$\bigcirc$\\
(4)~Persamaan garis singgung yang benar adalah $y=12x+18$.&$\bigcirc$&$\bigcirc$\\\hline
\end{tabular}}
\end{soal}

\begin{soal}{32 \normalfont\small[PG Majemuk Kompleks]}
\gambar[0.4\textwidth]{images/5/5-12.png}
Layang-layang $ABCD$; $A(4,3)$, $C(20,1)$; diagonal $BD$ sumbu simetri.
Persamaan garis $BD$: $y=mx+c$. Tentukan kebenaran setiap pernyataan berikut!

\par\smallskip
{\renewcommand{\arraystretch}{1.4}%
\begin{tabular}{@{}p{0.76\linewidth}cc@{}}
\hline\textbf{Pernyataan}&\textbf{Benar}&\textbf{Salah}\\\hline
(1)~Gradien $AC = \dfrac{1-3}{20-4} = -\dfrac{1}{8}$; gradien $BD = 8$.&$\bigcirc$&$\bigcirc$\\
(2)~Titik tengah $AC = (12,2)$; garis $BD$: $y-2=8(x-12)$, yaitu $y=8x-94$.&$\bigcirc$&$\bigcirc$\\
(3)~Dari $y=8x-94$: $m=8$ dan $c=-94$.&$\bigcirc$&$\bigcirc$\\
(4)~$m+c = 8+(-94) = -84$.&$\bigcirc$&$\bigcirc$\\\hline
\end{tabular}}
\end{soal}

\begin{soal}{33 \normalfont\small[PG Majemuk Kompleks]}
$\displaystyle\int_0^2 5(3a+bx)\,dx=\int_0^2 8x^2(3a+bx)\,dx=40$.
Tentukan kebenaran setiap pernyataan berikut!

\par\smallskip
{\renewcommand{\arraystretch}{1.4}%
\begin{tabular}{@{}p{0.76\linewidth}cc@{}}
\hline\textbf{Pernyataan}&\textbf{Benar}&\textbf{Salah}\\\hline
(1)~Dari integral pertama: $5\!\left[3ax+bx^2\right]_0^2 = 5(6a+4b)=40$, sehingga $3a+2b=4$.&$\bigcirc$&$\bigcirc$\\
(2)~Integral pertama yang benar: $5(6a+2b)=40$, sehingga $3a+b=4$.&$\bigcirc$&$\bigcirc$\\
(3)~Dari integral kedua: $64a+32b=40$, sehingga $8a+4b=5$.&$\bigcirc$&$\bigcirc$\\
(4)~Dari sistem persamaan diperoleh $a-b=7$.&$\bigcirc$&$\bigcirc$\\\hline
\end{tabular}}
\end{soal}

\begin{soal}{34 \normalfont\small[PG Majemuk Kompleks]}
$P(24,7)$; $A(28,9)$ dirotasi $90^\circ$ berlawanan jarum jam berpusat di $P$,
lalu dicerminkan terhadap sumbu-$x$. Bayangan akhir $B(a,b)$.
Tentukan kebenaran setiap pernyataan berikut!

\par\smallskip
{\renewcommand{\arraystretch}{1.4}%
\begin{tabular}{@{}p{0.76\linewidth}cc@{}}
\hline\textbf{Pernyataan}&\textbf{Benar}&\textbf{Salah}\\\hline
(1)~$\overrightarrow{PA}=(4,2)$; setelah rotasi 90° CCW: $(4,2)\mapsto(-2,4)$.&$\bigcirc$&$\bigcirc$\\
(2)~Titik hasil rotasi: $P+(-2,4)=(26,11)$.&$\bigcirc$&$\bigcirc$\\
(3)~Titik hasil rotasi adalah $(22,11)$; refleksi sumbu-$x$ menghasilkan $(22,-11)$.&$\bigcirc$&$\bigcirc$\\
(4)~$a+b=22+(-11)=11$.&$\bigcirc$&$\bigcirc$\\\hline
\end{tabular}}
\end{soal}

\begin{soal}{35 \normalfont\small[PG Majemuk Kompleks]}
\gambar[0.42\textwidth]{images/5/5-13.png}
Trapesium sama kaki $ABCD$, $AB\parallel CD$, $|AB|<|CD|$, $|AD|=15$,
$\sin\angle ADC=\frac{4}{5}$, keliling $K=78$.
Tentukan kebenaran setiap pernyataan berikut!

\par\smallskip
{\renewcommand{\arraystretch}{1.4}%
\begin{tabular}{@{}p{0.76\linewidth}cc@{}}
\hline\textbf{Pernyataan}&\textbf{Benar}&\textbf{Salah}\\\hline
(1)~Tinggi $h=|AD|\sin\angle ADC=15\cdot\dfrac{4}{5}=12$.&$\bigcirc$&$\bigcirc$\\
(2)~Komponen horizontal kaki $=\sqrt{15^2-12^2}=9$; $|CD|-|AB|=18$.&$\bigcirc$&$\bigcirc$\\
(3)~$|AB|+|CD|=K-2|AD|=48$; dari $|CD|-|AB|=18$: $|CD|=33$, $|AB|=15$.&$\bigcirc$&$\bigcirc$\\
(4)~Luas trapesium $=\tfrac{1}{2}(|CD|-|AB|)\cdot h=\tfrac{1}{2}(18)(12)=108$.&$\bigcirc$&$\bigcirc$\\\hline
\end{tabular}}
\end{soal}

\begin{soal}{36 \normalfont\small[PG Majemuk Kompleks]}
Lingkaran $C$: $x^2-6x+y^2-2y+6=0$. Rotasi $R_{(3,-4),\,\pi/2}(C)$
menghasilkan lingkaran berpusat $(k,l)$.
Tentukan kebenaran setiap pernyataan berikut!

\par\smallskip
{\renewcommand{\arraystretch}{1.4}%
\begin{tabular}{@{}p{0.76\linewidth}cc@{}}
\hline\textbf{Pernyataan}&\textbf{Benar}&\textbf{Salah}\\\hline
(1)~Pusat lingkaran $C$ setelah dilengkapi kuadrat: $(3,1)$.&$\bigcirc$&$\bigcirc$\\
(2)~Vektor dari pusat rotasi $(3,-4)$ ke $(3,1)$ adalah $(0,5)$; setelah rotasi CCW $\tfrac\pi2$: $(0,5)\mapsto(-5,0)$.&$\bigcirc$&$\bigcirc$\\
(3)~Pusat baru: $(3,-4)+(-5,0)=(2,-4)$, sehingga $k=2$ dan $l=-4$.&$\bigcirc$&$\bigcirc$\\
(4)~Pusat baru yang benar: $k=-2$ dan $l=-4$.&$\bigcirc$&$\bigcirc$\\\hline
\end{tabular}}
\end{soal}

\begin{soal}{37 \normalfont\small[PG Majemuk Kompleks]}
$\displaystyle\int(3x^4-2x^3)\,dx=\dfrac{1}{60}(ax^b+cx^d)+C$.
Tentukan kebenaran setiap pernyataan berikut!

\par\smallskip
{\renewcommand{\arraystretch}{1.4}%
\begin{tabular}{@{}p{0.76\linewidth}cc@{}}
\hline\textbf{Pernyataan}&\textbf{Benar}&\textbf{Salah}\\\hline
(1)~$\displaystyle\int 3x^4\,dx=\dfrac{3}{5}x^5=\dfrac{36}{60}x^5$, sehingga $a=36$ dan $b=5$.&$\bigcirc$&$\bigcirc$\\
(2)~$\displaystyle\int -2x^3\,dx=-\dfrac{1}{2}x^4=-\dfrac{30}{60}x^4$, sehingga $c=-30$ dan $d=5$.&$\bigcirc$&$\bigcirc$\\
(3)~Nilai $a-b+c-d=36-5+(-30)-4=-3$.&$\bigcirc$&$\bigcirc$\\
(4)~Nilai $a+c=36+(-30)=6$.&$\bigcirc$&$\bigcirc$\\\hline
\end{tabular}}
\end{soal}

\begin{soal}{38 \normalfont\small[PG Biasa]}
\gambar[0.45\textwidth]{images/5/5-14.png}
Pernyataan yang benar tentang grafik $y=f(x)$ di atas adalah \ldots
\begin{pil}
  \item $\lim_{x\to-2}f(x)=2$
  \item $\lim_{x\to2^-}f(x)=0$
  \item $\lim_{x\to2^+}f(x)=0$
  \item $f(5)=2$
  \item $\lim_{x\to5}f(x)=3$
\end{pil}
\end{soal}

\begin{soal}{39 \normalfont\small[PG Majemuk Kompleks]}
$a$, $b$, $c$ bilangan bulat dengan $\dfrac{2^a\cdot3^b\cdot5^c}{12^a\cdot10^b}=20$.
Tentukan kebenaran setiap pernyataan berikut!

\par\smallskip
{\renewcommand{\arraystretch}{1.4}%
\begin{tabular}{@{}p{0.76\linewidth}cc@{}}
\hline\textbf{Pernyataan}&\textbf{Benar}&\textbf{Salah}\\\hline
(1)~$12^a=2^{2a}3^a$ dan $10^b=2^b5^b$.&$\bigcirc$&$\bigcirc$\\
(2)~Ruas kiri $=2^{-a-b}\cdot3^{b-a}\cdot5^{c-b}=20=2^2\cdot5^1$.&$\bigcirc$&$\bigcirc$\\
(3)~Sistem persamaan: $-a-b=2$, $b-a=0$, $c-b=1$; sehingga $a=-1$, $b=-1$, $c=0$.&$\bigcirc$&$\bigcirc$\\
(4)~$a+b+c=-1+(-1)+0=-1$.&$\bigcirc$&$\bigcirc$\\\hline
\end{tabular}}
\end{soal}

\begin{soal}{40 \normalfont\small[PG Majemuk Kompleks]}
$f(x)=\begin{cases}x+3, & x<0\\ x^2-4x+5, & 0\leq x\leq3\\ 2x-4, & x>3\end{cases}$\\[0.5em]
Banyak bilangan bulat $x$ dengan $-3\leq x\leq6$ yang memenuhi $f(f(x))=2$.
Tentukan kebenaran setiap pernyataan berikut!

\par\smallskip
{\renewcommand{\arraystretch}{1.4}%
\begin{tabular}{@{}p{0.76\linewidth}cc@{}}
\hline\textbf{Pernyataan}&\textbf{Benar}&\textbf{Salah}\\\hline
(1)~Untuk $t>3$: $2t-4=2\Rightarrow t=3$ valid sebagai solusi $f(t)=2$.&$\bigcirc$&$\bigcirc$\\
(2)~Solusi $f(t)=2$ yang valid: $t=-1$, $t=1$, dan $t=3$.&$\bigcirc$&$\bigcirc$\\
(3)~Nilai $f(x)$ untuk $x=-3,-2,-1,\ldots,6$: $0,1,2,5,2,1,2,4,6,8$.&$\bigcirc$&$\bigcirc$\\
(4)~Banyak bilangan bulat yang memenuhi $f(f(x))=2$ adalah 2.&$\bigcirc$&$\bigcirc$\\\hline
\end{tabular}}
\end{soal}

% ===================================================================
%  SOAL 41--50
% ===================================================================

\begin{soal}{41 \normalfont\small[PG Biasa]}
Domain fungsi $g(x)=\dfrac{\sqrt{(x^2-5x+6)(x-1)}}{x^2-4x+3}$ adalah \ldots
\begin{pil}
  \item $(-\infty,1)\cup[2,3)$
  \item $(1,2]\cup(3,\infty)$
  \item $[1,2]\cup[3,\infty)$
  \item $(1,3)\cup\{2\}$
  \item $(1,\infty)\cup\{3\}$
\end{pil}
\end{soal}

\begin{soal}{42 \normalfont\small[PG Majemuk Kompleks]}
$f(1)=2$, $f'(1)=-3$; $h(x)=x^2f\!\left(\dfrac{1}{x}\right)$.
Tentukan kebenaran setiap pernyataan berikut!

\par\smallskip
{\renewcommand{\arraystretch}{1.4}%
\begin{tabular}{@{}p{0.76\linewidth}cc@{}}
\hline\textbf{Pernyataan}&\textbf{Benar}&\textbf{Salah}\\\hline
(1)~$h'(x)=2xf\!\left(\tfrac1x\right)-f'\!\left(\tfrac1x\right)$ (aturan hasil kali dan rantai).&$\bigcirc$&$\bigcirc$\\
(2)~$h'(1)=2f(1)-f'(1)$.&$\bigcirc$&$\bigcirc$\\
(3)~$h'(1)=2(2)-(-3)=7$.&$\bigcirc$&$\bigcirc$\\
(4)~Nilai $h'(1)=2f(1)+f'(1)=2(2)+(-3)=1$.&$\bigcirc$&$\bigcirc$\\\hline
\end{tabular}}
\end{soal}

\begin{soal}{43 \normalfont\small[PG Biasa]}
$f(x)=\begin{cases}\dfrac{x^3-1}{x-1}, & x<1\\ k, & x=1\\ ax+b, & x>1\end{cases}$\\[0.5em]
$f$ kontinu di $x=1$ dan turunan kiri di $x=1$ sama dengan dua kali turunan
kanan. Nilai $k+a+b=\ldots$
\begin{pilB}
  \item $4$ \item $5$ \item $6$ \item $7$ \item $8$
\end{pilB}
\end{soal}

\begin{soal}{44 \normalfont\small[PG Majemuk Kompleks]}
$H(x)=\displaystyle\int_0^{x^2+1}g(t)\,dt$; $g(2)=5$; $g'(2)=-1$.
Tentukan kebenaran setiap pernyataan berikut!

\par\smallskip
{\renewcommand{\arraystretch}{1.4}%
\begin{tabular}{@{}p{0.76\linewidth}cc@{}}
\hline\textbf{Pernyataan}&\textbf{Benar}&\textbf{Salah}\\\hline
(1)~$H'(x)=g(x^2+1)\cdot2x$ (Teorema Dasar Kalkulus + aturan rantai).&$\bigcirc$&$\bigcirc$\\
(2)~$H''(x)=2g(x^2+1)+2x\cdot g'(x^2+1)$.&$\bigcirc$&$\bigcirc$\\
(3)~$H''(1)=2g(2)+4g'(2)$.&$\bigcirc$&$\bigcirc$\\
(4)~$H''(1)=2(5)+4(-1)=6$.&$\bigcirc$&$\bigcirc$\\\hline
\end{tabular}}
\end{soal}

\begin{soal}{45 \normalfont\small[PG Majemuk Kompleks]}
$\sqrt{23-6\sqrt{10}}=\sqrt{a}-\sqrt{b}$; $a,b$ bulat positif; $a>b$.
Tentukan kebenaran setiap pernyataan berikut!

\par\smallskip
{\renewcommand{\arraystretch}{1.4}%
\begin{tabular}{@{}p{0.76\linewidth}cc@{}}
\hline\textbf{Pernyataan}&\textbf{Benar}&\textbf{Salah}\\\hline
(1)~$(\sqrt{a}-\sqrt{b})^2=a+b-2\sqrt{ab}$; membandingkan: $a+b=23$ dan $ab=90$.&$\bigcirc$&$\bigcirc$\\
(2)~Akar $t^2-23t+90=0$ adalah $t=15$ dan $t=6$, sehingga $a=15$, $b=6$.&$\bigcirc$&$\bigcirc$\\
(3)~Akar yang benar: $t=18$ dan $t=5$, sehingga $a=18$ dan $b=5$.&$\bigcirc$&$\bigcirc$\\
(4)~Nilai $a+b=23$.&$\bigcirc$&$\bigcirc$\\\hline
\end{tabular}}
\end{soal}

\begin{soal}{46 \normalfont\small[PG Majemuk Kompleks]}
$\sin(\pi x)=\cos(2\pi x)$; $0<x<\dfrac{1}{2}$. Tentukan kebenaran
setiap pernyataan berikut!

\par\smallskip
{\renewcommand{\arraystretch}{1.4}%
\begin{tabular}{@{}p{0.76\linewidth}cc@{}}
\hline\textbf{Pernyataan}&\textbf{Benar}&\textbf{Salah}\\\hline
(1)~Misalkan $\theta=\pi x$; persamaan menjadi $2\sin^2\theta+\sin\theta-1=0$.&$\bigcirc$&$\bigcirc$\\
(2)~$(2\sin\theta-1)(\sin\theta+1)=0$; untuk $0<\theta<\tfrac\pi2$: $\sin\theta=\tfrac12$.&$\bigcirc$&$\bigcirc$\\
(3)~$\theta=\dfrac{\pi}{6}$, sehingga $x=\dfrac{1}{6}$.&$\bigcirc$&$\bigcirc$\\
(4)~$\cos\!\left(\dfrac{\pi}{4}-\pi x\right)=\cos\dfrac{\pi}{12}=\dfrac{\sqrt{6}+\sqrt{2}}{4}$.&$\bigcirc$&$\bigcirc$\\\hline
\end{tabular}}
\end{soal}

\begin{soal}{47 \normalfont\small[PG Majemuk Kompleks]}
Kurva parametrik $x=t^2+t$, $y=t^3-3t$. Tentukan kebenaran pernyataan
tentang $\dfrac{d^2y}{dx^2}$ saat $t=1$ berikut!

\par\smallskip
{\renewcommand{\arraystretch}{1.4}%
\begin{tabular}{@{}p{0.76\linewidth}cc@{}}
\hline\textbf{Pernyataan}&\textbf{Benar}&\textbf{Salah}\\\hline
(1)~$\dfrac{dx}{dt}=2t+1$ dan $\dfrac{dy}{dt}=3t^2-3$.&$\bigcirc$&$\bigcirc$\\
(2)~$\dfrac{dy}{dx}=\dfrac{3t^2-3}{2t+1}$.&$\bigcirc$&$\bigcirc$\\
(3)~$\dfrac{d^2y}{dx^2}=\dfrac{6(t^2+t+1)}{(2t+1)^3}$.&$\bigcirc$&$\bigcirc$\\
(4)~Nilai $\dfrac{d^2y}{dx^2}$ saat $t=1$ adalah $\dfrac{3}{4}$.&$\bigcirc$&$\bigcirc$\\\hline
\end{tabular}}
\end{soal}

\begin{soal}{48 \normalfont\small[PG Majemuk Kompleks]}
Sistem $\begin{cases}x^2+2xy=3\\ x+y=a\end{cases}$ punya dua solusi real berbeda untuk $x$.
Tentukan kebenaran setiap pernyataan berikut!

\par\smallskip
{\renewcommand{\arraystretch}{1.4}%
\begin{tabular}{@{}p{0.76\linewidth}cc@{}}
\hline\textbf{Pernyataan}&\textbf{Benar}&\textbf{Salah}\\\hline
(1)~Substitusi $y=a-x$ menghasilkan $x^2+2ax-3=0$.&$\bigcirc$&$\bigcirc$\\
(2)~Agar dua solusi real berbeda: $\Delta=4a^2-12>0$.&$\bigcirc$&$\bigcirc$\\
(3)~$a^2>3$, yaitu $|a|>\sqrt{3}$.&$\bigcirc$&$\bigcirc$\\
(4)~Himpunan nilai $a$: $a<-\sqrt{3}$ atau $a>\sqrt{3}$.&$\bigcirc$&$\bigcirc$\\\hline
\end{tabular}}
\end{soal}

\begin{soal}{49 \normalfont\small[PG Majemuk Kompleks]}
$\displaystyle\lim_{x\to1}\dfrac{\sqrt{x^2+\frac{1}{x}}-\sqrt{ax+b}}{x-1}=L$ (berhingga).
Tentukan kebenaran setiap pernyataan berikut!

\par\smallskip
{\renewcommand{\arraystretch}{1.4}%
\begin{tabular}{@{}p{0.76\linewidth}cc@{}}
\hline\textbf{Pernyataan}&\textbf{Benar}&\textbf{Salah}\\\hline
(1)~Agar limit berhingga: pembilang $\to0$ saat $x\to1$, sehingga $a+b=2$.&$\bigcirc$&$\bigcirc$\\
(2)~$L=A'(1)-B'(1)$ dengan $A(x)=\sqrt{x^2+\frac1x}$, $B(x)=\sqrt{ax+b}$; hasilnya $L=\dfrac{1-a}{2\sqrt2}$.&$\bigcirc$&$\bigcirc$\\
(3)~Nilai $\displaystyle\lim_{x\to1}\dfrac{\sqrt{x^4+\frac{1}{x^2}}-\sqrt{ax^2+b}}{x-1}=2L$.&$\bigcirc$&$\bigcirc$\\
(4)~Nilai $\displaystyle\lim_{x\to1}\dfrac{\sqrt{x^4+\frac{1}{x^2}}-\sqrt{ax^2+b}}{x-1}=L$.&$\bigcirc$&$\bigcirc$\\\hline
\end{tabular}}
\end{soal}

\begin{soal}{50 \normalfont\small[PG Biasa]}
Banyak bilangan bulat $k$ sehingga $p_k(x)=x^3-3x+k$ mempunyai tiga akar
real berlainan adalah \ldots
\begin{pilB}
  \item $1$ \item $2$ \item $3$ \item $4$ \item $5$
\end{pilB}
\end{soal}

\end{document}
